\subsection{Contact orders in a perfect-residue specialization}
\label{subsec:boundary-specialization-contact}

This subsection gives a necessary contact condition for a specialization
that realizes an unequal-character comparison at the first unresolved
index.  A separate elementary example explains how an inseparable
residue extension can become ramification along a transverse curve.

\begin{proposition}[A necessary upstairs contact order]
\label{prop:boundary-specialization-contact}
Let $K$ be a complete discretely valued field of mixed characteristic
$(0,5)$ with algebraically closed residue field, $\mu _5\subset K$, and
$e_K=260$.  Suppose that $G,H$ belong to the standing height-two
filtration class, that there is an actual isomorphism
$G[5^4]\simeq H[5^4]$ over $K$, and that their special characters differ.
Let $F/K$ be a finite extension over which all the geometric
$5^4$-torsion points of $G$ are rational.

Put $E_*=624\cdot5^6$.  Suppose a local ring $S$ admits a local
homomorphism $S\to\mathcal O_F$ and elements $\pi,h\in S$ satisfying
\[
  5=\pi^{E_*}h,\qquad h\in S^\times,
  \qquad q:=v_F(\pi)>0,
\]
where $v_F$ is normalized to have value group $\mathbf Z$.
Then $5\mid q$.  In particular, the displayed chart cannot realize such
a comparison by an upstairs transverse curve with $q=1$.
\end{proposition}

\begin{proof}
The hypotheses are the case $p=5$, $r=1$ of
Theorem~\ref{thm:first-boundary-normal-form}.  Since the residue field is
already algebraically closed, Corollary~\ref{cor:boundary-point-field-degrees}
applies after a common tame good extension $L/K$, which may be chosen
Galois of degree dividing $24$.  If $Q$ is an exact $5^4$-torsion point
whose $5^3$-multiple is a nonzero canonical point, that corollary gives
\[
  [L(Q):L]=5^6.
\]
The full-level hypothesis on $F$ implies $Q\in G(FL)$.  Consequently
$5^6\mid[FL:L]$.  The two tower formulas give
\[
  [FL:L][L:K]=[FL:F][F:K].
\]
Both $[L:K]$ and $[FL:F]$ are prime to $5$, because $L/K$ is Galois
and tame.  It follows that
\[
  5^6\mid[F:K].
\]
This argument requires neither $F/K$ nor a whole point-field tower
to be Galois.

Every finite extension of $K$ has residue field equal to the given
algebraically closed field.  The degree formula for finite extensions
of complete discrete valuation fields therefore gives
$[F:K]=e(F/K)$; see
\citet[Section~15.113, Remark~15.113.6]{StacksDVR}.
The local homomorphism carries $h$ to a unit, so
\[
  v_F(5)=E_*q,
  \qquad
  [F:K]=\frac{v_F(5)}{e_K}
       =\frac{624\cdot5^6q}{260}
       =12\cdot5^5q.
\]
The divisibility $5^6\mid[F:K]$ now implies $5\mid q$.
\end{proof}

Here $\pi$ is a local equation of the upstairs divisor, and the unit
hypothesis requires that $5/\pi^{E_*}$ remain invertible at the chosen
closed point.  The proposition does not apply in this form if that
factor vanishes.  It also does not exclude a downstairs transverse
curve whose lift has contact $5$ or a larger multiple of $5$ upstairs.
For example, if a finite extension of degree $12\cdot5^6$ were retained
under these hypotheses, the displayed degree identity would force
$q=5$.  No existence of such an extension or of the required comparison
is asserted.

The specialization must be made on a model.  There is no $k$-algebra
homomorphism $k(t)\to k$: if $t$ mapped to $c\in k$, the invertible
element $t-c$ would map to zero.  Thus one must choose the closed point
and transverse curve before inverting all nonzero residue functions.

\begin{example}[An inseparable residue extension becoming ramification]
\label{ex:fierce-transverse-specialization}
Let $k=\overline{\mathbf F}_5$.  Choose a complete discrete valuation
field $K_0$ with residue field $k$, $\mu _5\subset K_0$, and
$e_{K_0}=260$.  Such a field is obtained from $\operatorname{Frac}W(k)$
by adjoining a primitive fifth root of unity $\zeta$, then a root
$\varpi$ of
\[
  \varpi^{65}=\zeta-1.
\]
The translated cyclotomic polynomial is Eisenstein of degree $4$ over
$W(k)$, and the second polynomial is Eisenstein of degree $65$ over
the resulting valuation ring.  Hence $\varpi$ is a uniformizer of
$K_0$ and $v_{K_0}(5)=4\cdot65=260$.

Consider the finite flat degree-five cover
\[
 R=\mathcal O_{K_0}[s]_{(\varpi,s-1)},
 \qquad B=R[a]/(a^5-s).
\]
At the generic point of the divisor $\varpi=0$, complete the downstairs
discrete valuation ring and denote its fraction field by $E$.
Its residue field is $k(s)$.  The polynomial $T^5-s$ is irreducible
over $k(s)$, so the corresponding integral algebra over the complete
valuation ring of $E$ has residue field $k(s^{1/5})$.  It is local,
with maximal ideal generated by $\varpi$, and is therefore a discrete
valuation ring.  It is free of rank $5$ over the downstairs ring.
Consequently the extension $E(a)/E$ has
\[
  e(E(a)/E)=1,\qquad
  f(E(a)/E)=5,
\]
with purely inseparable residue extension.  It is a cyclic extension
of fields because $\mu _5\subset E$.

Now restrict the cover to the downstairs curve $s=1+\varpi$.
This curve is transverse to $\varpi=0$, since its coordinate ring is
$\mathcal O_{K_0}$ and its intersection with the divisor has length
one.  Putting $z=a-1$ gives
\[
  z^5+5z^4+10z^3+10z^2+5z-\varpi=0.
\]
This polynomial is Eisenstein at $\varpi$.  The specialized fraction
field extension therefore has degree $5$, ramification index $5$,
and residue degree $1$.  The downstairs absolute index remains
$260$, while the upstairs contact order with $\varpi=0$ is $5$.

By comparison, the downstairs transverse curve $s=1$ gives $a^5=1$;
its generic fibre splits into five copies of $K_0$.
Thus the generic inseparable residue extension determines neither
retention nor loss of degree under transverse specialization.
This example concerns a Kummer cover only.  It supplies neither a
Drinfeld-level specialization nor a pair of groups with unequal
special characters and isomorphic $5^4$ torsion.
\end{example}
